\scanpage{154}
\deleted{fait $\lambda=p\lambda^+$}], et les hom (17) $P\xrightarrow{d_P}E\otimes\Omega_{S/T}$ satisfaisant les conditions (18), en associant à $\lambda^+$ sa réduction dans $P^1_{S/T}(E)^+=E\otimes\Omega_{S/T}$.

\uline{NB} Prenant $P$ trivial sur $E$, ceci prouve que l'hom \deleted{de réduction} de réduction $P^\infty_{S/T}(E)^+\longrightarrow P^1_{S/T}(E)^+$ induit un \deleted{\unsure{biject|biject.|bijection}} \uline{isomorphisme de faisceaux}
\[
 P^\infty_{S/T}(E)^{+\natural}\xrightarrow{\sim}{\unsure{Z|\mathbb Z}}^{1}(E\otimes\Omega^\bullet_{S/T})\,,
\]
en prenant au deuxième m\textsuperscript{b} comme stratification sur $P^\infty_{S/T}(E)^+$ celle déduite de la strat. can. de $P^\infty_{S/T}(E)/\delta E$ par transport de structure à l'aide de (29). Ainsi, \uline{une} 1-\uline{forme fermée sur} $S/T$ \inserted{à \uline{coeff. dans} $E$} \uline{définit un jet d'ordre} $\infty$ \inserted{$\lambda^+$} \uline{de} $E$ \uline{sur} $S/T$, satisfaisant la condition (28). A la forme exacte $d^0_E(\unsure{\xi|\zeta})$ est associé ainsi \deleted{\unsure{$\lambda(\xi)=d_E(\xi)$}} \deleted{\unsure{$\lambda^+(\xi)$}}, $\lambda(\xi)=d^\infty_E(\xi)$ i.e. $\lambda^+(\xi)=d^\infty_E(\xi)-\delta^\infty_E(\xi)=d^+_E(\xi)^{\unsure{+|1}}$.

4) \uline{Equations de Picard Fuchs}.

Supposons encore qu'on ait $E$ sur $S$, avec stratification $\delta^\infty_E=\delta$. Supposons donné de plus une suite exacte de modules qu. coh.
\[
 0\longrightarrow E'\mathrel{\overset{u}{\text{---}}}E\mathrel{\overset{v}{\text{---}}}E''\mathrel{\text{---}}0 . \tag*{(\unsure{33|38})}
\]
On ne suppose pas que $E'$ soit un sous-module stratifié. \deleted{\unsure{Supposons}} Supposons donné un
% archive sheet number: 153
